\chapter{Einleitung}
Im Rahmen des Projektmoduls entwickelten die Studierenden in Kleingruppen eigene Projektideen mit Bezug zum Oldenburger Computer-Museum (OCM). Ziel des Moduls war es, theoretische Kenntnisse aus dem Studium in einem praxisnahen Umfeld anzuwenden und eigenständig Lösungen für reale Anwendungsfälle zu konzipieren und umzusetzen.

Zu Beginn des Semesters fanden sich die Studierenden selbstständig zu Gruppen zusammen und entwickelten eine Projektidee. Anschließend wurden die Konzepte weiter ausgearbeitet und während des Semesters schrittweise umgesetzt. Die Projektarbeit erfolgte dabei weitgehend eigenverantwortlich. Im Rahmen regelmäßiger Treffen in der Hochschule stellten die Gruppen ihren aktuellen Arbeitsstand vor, berichteten über Fortschritte und Herausforderungen und erhielten Feedback zum weiteren Vorgehen.

Während der gesamten Projektlaufzeit standen die Gruppen in engem Austausch mit dem betreuenden Dozenten, Prof. Dr. Paelke. Außerdem gab es regelmäßigen Kontakt mit der Ansprechperson des Oldenburger Computer-Museums, um fachliche Fragen zu klären und die entwickelten Lösungen an die Anforderungen des Museums anzupassen.

Die vorliegende Dokumentation fasst die Ergebnisse der einzelnen Projektgruppen zusammen. Jedes Kapitel beschreibt ein eigenständiges Projekt und gibt einen Überblick über die Ausgangssituation, den Entwicklungsprozess sowie die erzielten Ergebnisse.

